\section{端云协同投机解码方案设计}

\subsection{系统总体设计}

本文系统实现了一个跨设备投机解码流程：Android 端在本地运行草稿模型，根据当前已接受前缀生成候选词元；桌面端运行目标验证服务，对候选词元执行接受、拒绝和修正判定；客户端依据验证结果更新本地草稿状态，并继续下一轮提案。系统通过 HTTP 和 JSON 消息进行交互，围绕同一投机会话维护草稿侧与目标侧状态。

本文不展开 Android 本地推理和普通远程推理的系统设计，而将二者作为投机解码系统的实验前提：前者用于确认移动端具备真实模型运行和草稿生成能力，后者用于确认移动端与桌面端之间具备基础通信和远程推理条件。

图 \ref{fig:system-architecture} 给出了本文系统总体架构。
\begin{figure}[htbp]
  \centering
  \begin{tikzpicture}[
    scale=0.86,
    transform shape,
    node distance=10mm and 9mm,
    box/.style={draw, rounded corners, align=center, minimum width=26mm, minimum height=10mm},
    line/.style={-Latex, thick}
  ]
    \node[box] (user) {用户输入};
    \node[box, right=of user] (android) {Android 客户端};
    \node[box, below left=of android] (draft) {草稿会话\\本地草稿生成};
    \node[box, below=of android] (protocol) {协议层\\start / propose\\fallback / close};
    \node[box, right=of protocol] (desktop) {Desktop\\Inference Service};
    \node[box, right=of desktop] (target) {目标会话\\目标验证状态};
    \node[box, above=of target] (runtime) {llama.cpp / llama-server\\目标模型运行时};
    \node[box, below=of desktop] (log) {实验数据采集};
    \node[box, above=of android] (output) {生成结果展示};

    \draw[line] (user) -- (android);
    \draw[line] (android) -- (draft);
    \draw[line] (android) -- (protocol);
    \draw[line] (protocol) -- (desktop);
    \draw[line] (desktop) -- (target);
    \draw[line] (target) -- (runtime);
    \draw[line] (runtime) -- (target);
    \draw[line] (desktop) -- (log);
    \draw[line] (android) -- (output);
  \end{tikzpicture}
  \caption{系统总体架构图}
  \label{fig:system-architecture}
\end{figure}

从模块角度看，系统可分为四部分：草稿生成模块、目标验证模块、协议与会话管理模块、性能统计模块。草稿生成模块位于 Android 端，负责读取已接受前缀、维护草稿运行时状态并生成候选词元或候选树。目标验证模块位于桌面端，负责在目标模型上下文中检查候选，给出可接受长度和修正词元。协议与会话管理模块连接两端运行时，保证 \texttt{start}、\texttt{propose}、\texttt{fallback} 和 \texttt{close} 等消息在同一会话语义下工作。性能统计模块记录每轮提案规模、接受数量、草稿侧耗时和验证侧耗时，为后续实验分析提供数据。各模块功能如表 \ref{tab:modules} 所示。
\begin{table}[htbp]
  \centering
  \caption{系统主要模块及职责}
  \label{tab:modules}
  \small
  \begin{tabularx}{0.98\textwidth}{p{24mm}p{22mm}X}
    \toprule
    模块 & 部署位置 & 主要职责 \\
    \midrule
    草稿生成模块 & Android 端 & 维护草稿会话，生成候选词元或候选树 \\
    目标验证模块 & 桌面端 & 维护目标会话，执行接受、拒绝与修正判定 \\
    协议管理模块 & Android + 桌面端 & 管理 \texttt{start}、\texttt{propose}、\texttt{fallback}、\texttt{close} 生命周期 \\
    性能统计模块 & 桌面端为主 & 采集验证耗时、接受数量、提案规模和词元统计信息 \\
    \bottomrule
  \end{tabularx}
\end{table}

\subsection{投机解码流程设计}

本文系统的投机解码流程包括五个关键阶段：会话建立、草稿生成、目标验证、状态推进和会话结束。

第一阶段为会话建立。用户在 Android 端输入提示词后，客户端向桌面端发送 \texttt{start} 请求。桌面端据此创建投机会话，并为其分配独立的目标会话，用于记录目标侧验证状态。

第二阶段为草稿生成。Android 端基于当前已接受前缀启动或恢复草稿会话，生成一段线性提案或一棵浅层草稿树。

第三阶段为目标验证。Android 端通过 \texttt{propose} 请求发送候选词元或草稿树，桌面端在目标会话中读取当前状态后执行验证，并返回接受数量、拒绝位置和修正词元。

第四阶段为状态推进。若提案中的前若干词元被接受，则已接受前缀向前推进；若发生拒绝，则系统只保留已接受前缀与修正词元所形成的新权威前缀，其他未被接受的草稿部分全部丢弃。

第五阶段为继续迭代或关闭会话。只要未达到结束条件，系统便围绕新的已接受前缀继续下一轮提案与验证；若会话终止，则发送 \texttt{close} 请求释放资源。

图 \ref{fig:sequence} 给出了端云协同投机解码的时序过程。
\begin{figure}[htbp]
  \centering
  \begin{tikzpicture}[
    scale=0.86,
    transform shape,
    lifeline/.style={draw, minimum width=22mm, minimum height=8mm, align=center},
    msg/.style={-Latex, thick},
    note/.style={font=\small}
  ]
    \node[lifeline] (u) at (0,0) {用户};
    \node[lifeline] (a) at (3.2,0) {Android 端};
    \node[lifeline] (d) at (6.4,0) {桌面端服务};
    \node[lifeline] (t) at (9.6,0) {目标运行时};

    \foreach \x in {u,a,d,t} {
      \draw[dashed] (\x.south) -- ++(0,-8.2);
    }

    \draw[msg] (u) -- node[above,note]{输入提示} (a);
    \draw[msg] ($(a.south)+(0,-0.6)$) -- node[above,note]{start} ($(d.south)+(0,-0.6)$);
    \draw[msg] ($(d.south)+(0,-1.2)$) -- node[above,note]{初始化目标会话} ($(t.south)+(0,-1.2)$);
    \draw[msg] ($(t.south)+(0,-1.8)$) -- node[above,note]{会话就绪} ($(d.south)+(0,-1.8)$);
    \draw[msg] ($(d.south)+(0,-2.4)$) -- node[above,note]{start response} ($(a.south)+(0,-2.4)$);

    \draw[msg] ($(a.south)+(0,-3.3)$) -- node[above,note]{生成草稿提案} ++(0,-0.7);
    \draw[msg] ($(a.south)+(0,-4.4)$) -- node[above,note]{propose} ($(d.south)+(0,-4.4)$);
    \draw[msg] ($(d.south)+(0,-5.0)$) -- node[above,note]{基于已接受前缀验证} ($(t.south)+(0,-5.0)$);
    \draw[msg] ($(t.south)+(0,-5.6)$) -- node[above,note]{返回接受数量与修正词元} ($(d.south)+(0,-5.6)$);
    \draw[msg] ($(d.south)+(0,-6.2)$) -- node[above,note]{返回验证结果} ($(a.south)+(0,-6.2)$);
    \draw[msg] ($(a.south)+(0,-6.8)$) -- node[above,note]{应用结果并推进状态} ++(0,-0.7);

    \draw[msg] ($(a.south)+(0,-8.0)$) -- node[above,note]{close} ($(d.south)+(0,-8.0)$);

    \node[note, anchor=west] at (10.7,-0.75) {阶段一：会话建立};
    \node[note, anchor=west] at (10.7,-3.55) {阶段二：草稿生成};
    \node[note, anchor=west] at (10.7,-4.95) {阶段三：目标验证};
    \node[note, anchor=west] at (10.7,-6.65) {阶段四：状态推进};
    \node[note, anchor=west] at (10.7,-8.15) {阶段五：会话结束};
  \end{tikzpicture}
  \caption{投机解码五阶段时序图}
  \label{fig:sequence}
\end{figure}

图 \ref{fig:sequence} 中的时序可对应到五个阶段：\texttt{start} 请求和目标会话初始化对应会话建立阶段；Android 端本地生成候选对应草稿生成阶段；\texttt{propose} 请求和目标运行时检查对应目标验证阶段；验证结果返回和 Android 端应用结果对应状态推进阶段；最后的 \texttt{close} 请求对应会话结束阶段。通过这五个阶段，系统把一次完整生成任务拆分为多轮“提案--验证--推进”循环。

\subsection{协议与会话管理设计}

为了支持多轮投机会话，系统定义了 \texttt{start}、\texttt{propose}、\texttt{fallback} 和 \texttt{close} 四类基础消息。其中，\texttt{start} 用于建立会话并初始化目标侧状态；\texttt{propose} 用于提交草稿候选并获取验证结果；\texttt{fallback} 用于在协同推理不可用或验证失败时退回普通远程推理；\texttt{close} 用于释放投机会话和目标会话。

协议设计包含三项要求。第一，消息语义保持稳定，不同验证器实现均围绕接受数量、拒绝位置和修正词元这组语义工作。第二，协议保留必要的性能统计字段，用于量化协同开销和接受行为。第三，协议同时支持线性提案、树形提案和真实词元表示。

系统消息结构可归纳如表 \ref{tab:protocol} 所示。
\begin{table}[htbp]
  \centering
  \caption{投机解码协议消息设计}
  \label{tab:protocol}
  \small
  \begin{tabularx}{0.98\textwidth}{p{18mm}p{18mm}p{40mm}X}
    \toprule
    消息类型 & 发送方 & 主要字段 & 作用 \\
    \midrule
    \texttt{start} & Android 端 & \texttt{requestId}、\texttt{systemPrompt}、\texttt{userPrompt}、\texttt{model} & 创建投机会话和目标会话 \\
    \texttt{propose} & Android 端 & \texttt{sessionId}、\texttt{proposedTokenIds}、\texttt{draftTree}、\texttt{draftStep} & 提交草稿候选并请求验证 \\
    \texttt{fallback} & Android 端 & \texttt{sessionId}、普通生成参数 & 退回普通远程推理方案 \\
    \texttt{close} & Android 端 & \texttt{sessionId} & 关闭会话并释放资源 \\
    \bottomrule
  \end{tabularx}
\end{table}

除消息本身外，会话状态也是系统设计重点。本文将状态显式区分为草稿会话与目标会话。前者保存草稿运行时、已接受前缀和草稿配置；后者保存目标验证状态、已接受文本、词元进度以及验证连续性相关字段。两类状态分别对应草稿侧与目标侧的状态归属。

\subsection{关键问题分析}

在端云协同投机解码中，本文重点面对四类问题。

第一是状态同步问题。草稿侧和目标侧运行在不同设备上，一旦任意一轮已接受前缀没有被双方一致应用，后续提案和验证就会失去共同前缀基础，导致接受率迅速下降或会话失效。系统通过草稿会话和目标会话分别保存权威前缀，并在每轮验证后只提交目标侧确认过的词元，以保证双方围绕同一上下文继续生成。

第二是验证真实性问题。投机解码的收益建立在目标模型真实验证之上，若目标侧只进行文本代理比较或使用不一致的上下文，接受结果便不能代表目标模型分布。系统在桌面端目标会话中基于真实词元 id 执行验证，并返回可用于推进会话状态的接受数量、拒绝位置和修正词元。

第三是词元空间一致性问题。若 Android 草稿词元和桌面端目标词元不处于统一的真实 id 空间中，则标准接受规则难以稳定实现\cite{li2024eagle}。系统要求草稿提案、目标验证和修正词元均使用真实词元 id 表示，避免在跨设备传输中引入文本投影或编码映射误差。

第四是中间协同开销问题。投机解码理论上减少目标模型逐词元解码次数，但在端云场景中，每一轮 \texttt{propose} 都伴随请求处理、状态同步和验证器运行时开销。如果这些中间成本过高，则端到端收益会被削弱。系统通过持续化会话、减少请求字段和记录分项耗时来控制并量化该问题，使实验能够区分草稿侧、验证侧和跨设备协同边界的开销来源。
